\section{Summary}\label{sec:ch3_summary}

This chapter has established the formal and conceptual foundations for the remainder of the dissertation. We defined open-world environments as settings in which an agent's knowledge is incomplete with respect to the true state of the world, and formalized novelty as a mismatch between the agent's internal models and its environment. The taxonomy of novelty types---prohibitive, obstructive, beneficial, and nuisance---provides a principled vocabulary for reasoning about how novelty affects task performance, while the classification by domain element---object, attribute, and action novelty---characterizes what aspects of the environment have changed.

We then developed the agent framework, integrating symbolic planning (\cref{sec:symbolic_planning}), reinforcement learning (\cref{sec:rl}), and the integrated planning task (IPT) formalism (\cref{sec:hybrid}) that couples them through a detector $\Abstraction : \States \rightarrow \SymStates$ and an executor mapping $e : \Operators \rightarrow \Executors$. The framework distinguishes planning success (plannability) from execution success (solvability), making explicit where novelty can manifest: at the symbolic level through incomplete models, or at the execution level through failed executors while symbolic plans remain valid.

Finally, we introduced the concept of structured abstractions and argued that the choice of abstraction determines whether novelty-induced mismatch leads to brittle failure or can be managed through efficient adaptation and generalizable behavior. Action abstractions expose the compositional structure of plans, enabling failure localization and targeted learning. Interaction abstractions capture the physically meaningful invariants of manipulation tasks, enabling transfer across objects and embodiments.

The subsequent chapters instantiate this framework. \Cref{ch:perspectives} situates the dissertation within the broader literature on open-world planning, learning, execution recovery, and agents built with foundation models.
\Cref{ch:benchmarks} introduces benchmark environments and evaluation methodologies that operationalize the novelty taxonomy through controlled injection of object, attribute, and action novelties, and provide metrics for measuring adaptation and recovery. \Cref{ch:rapid_learn} develops hybrid planning--learning methods that exploit action abstractions to achieve efficient adaptation under novelty. \Cref{ch:generalization} presents force-based learning approaches that use interaction abstractions to enable generalizable robotic manipulation. In each case, chapter-specific related work is discussed in the context of the particular contribution.